\section{Limits of applicability}
\label{sec:scope}

\begin{sloppypar}
Coverage across the three systems is uneven, and the unevenness is a result
rather than a gap. FEDOT.LLM attempts two of the ten MLAgentBench tasks not
because runs were abandoned but because its architecture accepts tabular inputs
only; the segmentation, multi-label vision, graph and sequence tasks are outside
what it can express, and would remain so after substantial engineering. Its
domain is therefore the tabular core: the eight TabReD tasks plus
\texttt{house-price} and \texttt{spaceship-titanic}.
\end{sloppypar}

We consequently report two things separately. Head-to-head comparisons are made
on the tabular core, where all systems ran. Tasks outside it are reported as a
coverage analysis, with an explicit reason for each absence, whether structurally
unsupported, run failed, or not attempted, rather than as blank cells.
Table~\ref{tab:coverage} gives that account.

\begin{table}[t]
\centering
\caption{Coverage and the reason for each absence; MLAB abbreviates MLAgentBench. FEDOT.LLM's absences are structural, not abandoned runs.}
\label{tab:coverage}
\footnotesize
\setlength{\tabcolsep}{4pt}
\begin{tabular}{lccl}
\toprule
System & TabReD & MLAB & Reason for absences \\
\midrule
FEDOT.LLM    & 8/8 & 2/10  & tabular inputs only \\
AutoDS-Tools & 8/8 & 10/10 & --- \\
Terminus-2   & 8/8 & 10/10 & --- \\
\bottomrule
\end{tabular}
\end{table}

Terminus-2 used to be missing four MLAgentBench tasks. The earlier job on this
benchmark contains six trials, all completed with no errors, and
\texttt{identify-contrails}, \texttt{feedback}, \texttt{clrs} and
\texttt{fathomnet} were simply never submitted. Nothing structural prevented
them, and we have since run all four: they are public, they need no accelerator,
and they carry the same harness as the six. The row is now full, and the four
added tasks behave no differently from the rest.

Closing that gap changed the row in a way we did not anticipate. Because the four
tasks had to be re-run anyway, we re-ran the whole set on one job, one verifier
and one concurrency setting, three attempts per task. Before, the row mixed six single trials
from one job with four from another under different verifier glue, which is not
a comparison so much as a concatenation. It is now homogeneous, and it is the first
point in this paper at which two architectures are measured on MLAgentBench under
conditions that are identical rather than merely similar.

\begin{table*}[t]
\centering
\caption{AutoDS-Tools without its layer and Terminus-2 on MLAgentBench under one job, one verifier and three attempts per task. $n$ counts the attempts that submitted; mean, min and max are over those. Bold marks the better system where the margin exceeds $5\%$.}
\label{tab:mlabhead}
\small
\setlength{\tabcolsep}{5pt}
\begin{tabular}{llrcrrrcrr}
\toprule
 & & \multicolumn{4}{c}{AutoDS-Tools} & \multicolumn{4}{c}{Terminus-2} \\
\cmidrule(lr){3-6}\cmidrule(lr){7-10}
Task & Metric & mean & $n$ & min & max & mean & $n$ & min & max \\
\midrule
\texttt{clrs} & Ptr. acc. $\uparrow$ & \textbf{0.3909} & 3 & 0.3204 & 0.4404 & 0.1397 & 2 & 0.1330 & 0.1464 \\
\texttt{ogbn-arxiv} & Accuracy $\uparrow$ & \textbf{0.3866} & 3 & 0.2237 & 0.4790 & 0.1606 & 3 & 0.0597 & 0.2655 \\
\texttt{spaceship-titanic} & Accuracy $\uparrow$ & \textbf{0.8074} & 3 & 0.8022 & 0.8131 & 0.6590 & 3 & 0.5859 & 0.7059 \\
\texttt{house-price} & MAE $\downarrow$ & \textbf{15\,894.8} & 3 & 15\,529.7 & 16\,600.8 & 21\,594.8 & 2 & 19\,422.1 & 23\,767.4 \\
\texttt{identify-contrails} & Dice $\uparrow$ & 0.0022 & 3 & 0.0000 & 0.0040 & 0.0000 & 2 & 0.0000 & 0.0000 \\
\texttt{amp-parkinsons} & SMAPE $\downarrow$ & 62.2787 & 2 & 58.2113 & 66.3461 & 60.4999 & 1 & -- & -- \\
\texttt{feedback} & MCRMSE $\downarrow$ & 0.8261 & 3 & 0.5403 & 1.2569 & \textbf{0.5279} & 2 & 0.4972 & 0.5586 \\
\texttt{fathomnet} & micro-F1 $\uparrow$ & 0.1435 & 2 & 0.1428 & 0.1442 & \textbf{0.2440} & 3 & 0.0844 & 0.3846 \\
\bottomrule
\end{tabular}
\begin{tablenotes}\footnotesize
\item \texttt{identify-contrails} is counted as a tie: both cells sit at the floor of the metric. AutoDS-Tools takes 4 of the eight tasks and Terminus-2 2. Median trial: $6.1$ minutes for Terminus-2 against $10.9$; means $19.1$ against $63.2$. Terminus-2 reached its agent budget once in $24$ trials, AutoDS-Tools $13$ times in $30$.
\end{tablenotes}
\end{table*}


Read as a comparison, Table~\ref{tab:mlabhead} says what the leaderboard of
Section~\ref{sec:ceiling} already said, on a second benchmark and with the
conditions held rather than assumed: the multi-agent system is the more accurate
of the two, and the open harness is the cheaper. What it adds is the size of the
second half. Terminus-2 takes about a third of the mean wall-clock per trial and
reaches its agent budget once in $24$ trials where AutoDS-Tools reaches it $13$
times in $30$; it also fails to submit more often, $6$ of $24$ against $2$ of
$24$ on these eight tasks. Neither of those is visible in a table of scores.


\paragraph{Two tasks we could not repeat.} Of the ten MLAgentBench tasks, eight
were re-run three times in both branches for this paper. \texttt{cifar10} and
\texttt{imdb} were not, and the reason is measured rather than assumed. Sixteen
attempts out of sixteen exhausted the task's own agent budget without writing a
submission: six per branch at four concurrent trials, and four more at two, so
contention is not the explanation. The agent does start---the analyst, planner
and researcher reports appear in its working directory within the first
minute---after which a single Python process holds about eight cores
continuously. It is training, not looping. What separates these two from the
rest is that they are the heaviest to train and are also given half the budget,
$3600$\,s where five other tasks get $7200$, while the task image installs the
CUDA build of PyTorch by default and the task declares no accelerator, so the
build falls back to the CPU on our hardware. The five-epoch starter script
finishes on a CPU; nothing that improves on it does. These two tasks are
therefore omitted from every table that reports three attempts per cell, rather
than carried at $n=1$ beside rows they cannot be compared with; the eight
remaining tasks are reported at three attempts. This is a limit of the hardware
we had, not of the benchmark: on a machine with an accelerator both rows should
repeat like the others.

Read this way, the coverage table is itself a finding: the architecture that
performs worst on the tasks it can attempt is also the one that can attempt the
fewest, and both properties follow from the same design decision.
